\documentclass[10pt]{article}
\usepackage[T1]{fontenc}
\usepackage[utf8]{inputenc}
\usepackage[french]{babel}
\usepackage{amsmath,amssymb,mathrsfs,mathpazo}
\usepackage{geometry}
\geometry{a4paper,margin=2cm}
\usepackage{graphicx,array,tabularx,tabularray,colortbl,multirow,multicol,booktabs,enumitem,pifont,ulem,fancybox,diagbox,cancel}
\usepackage{pstricks,pstricks-add,pst-plot,pst-tree,pst-node,pst-text,pst-func,pst-math,pst-bspline,pst-3dplot,pst-eucl,pst-solides3d}
\let\pscalculate\undefined
\usepackage{pst-bezier}
\usepackage{tikz,pgfplots,tkz-tab}
\usetikzlibrary{arrows,arrows.meta,patterns,calc,positioning,trees,decorations.markings,shapes.geometric,shapes.misc,shapes.arrows,intersections,decorations.pathmorphing,decorations.pathreplacing,backgrounds,shadows,babel}
\pgfplotsset{compat=1.18}
\usepgfplotslibrary{fillbetween}
\usepackage[np]{numprint}
\usepackage{listings,makecell,pas-tableur,twoopt}
\let\barre\undefined
\usepackage[pstricks]{tablvar}
\definecolor{darkgreen}{RGB}{0,100,0}
\providecommand{\red}{\color{red}}
\providecommand{\blue}{\color{blue}}
\providecommand{\green}{\color{green}}
\providecommand{\black}{\color{black}}
\providecommand{\white}{\color{white}}
\providecommand{\gray}{\color{gray}}
\providecommand{\cyan}{\color{cyan}}
\providecommand{\magenta}{\color{magenta}}
\providecommand{\yellow}{\color{yellow}}
\providecommand{\violet}{\color{violet}}
\providecommand{\orange}{\color{orange}}
\providecommand{\euro}{\texteuro}
\providecommand{\EUR}{\texteuro}
\providecommand{\ieme}{\textsuperscript{e}}
\providecommand{\degres}{\ensuremath{^\circ}}
\providecommand{\degree}{\ensuremath{^\circ}}
\providecommand{\vv}[1]{\overrightarrow{#1}}
\providecommand{\Coord}[1]{#1}
\providecommand{\decosix}{}
\providecommand{\starredbullet}{\ensuremath{\bullet}}
\pagestyle{empty}
\setlength{\parindent}{0pt}
\usepackage{adjustbox}\DeclareUnicodeCharacter{2236}{:}\DeclareUnicodeCharacter{2190}{\ensuremath{\leftarrow}}
\providecommand{\R}{\mathbb{R}}
\providecommand{\N}{\mathbb{N}}
\providecommand{\Z}{\mathbb{Z}}
\providecommand{\Q}{\mathbb{Q}}
\providecommand{\D}{\mathbb{D}}
\providecommand{\C}{\mathbb{C}}
\providecommand{\e}{{\mathrm{e}}}
\providecommand{\vect}[1]{\overrightarrow{#1}}
\providecommand{\vectt}[1]{\overrightarrow{\mathrm{#1}}}
\providecommand{\vvt}[1]{\overrightarrow{\mathrm{#1}}}
\providecommand{\np}[2][]{{#2~#1}}
\providecommand{\vv}[1]{{\overrightarrow{#1}}}
\renewcommand{\d}{{\mathrm{d}}}
\providecommand{\degres}{{^\circ}}
\providecommand{\degre}{{^\circ}}
\providecommand{\blue}{\color{blue}}
\providecommand{\red}{\color{red}}
\providecommand{\enskip}{\quad}
\providecommand{\llbracket}{\lbrack\!\lbrack}
\providecommand{\rrbracket}{\rbrack\!\rbrack}
\providecommand{\euro}{{\text{€}}}
\providecommand{\no}{{\text{n°}}}
\providecommand{\interval}[2]{{[#1;#2]}}
\providecommand{\intervFO}[2]{{[#1;#2[}}
\providecommand{\dotfill}{\ldots}

\providecommand{\ds}{\displaystyle}
\providecommand{\textding}[1]{#1}
\providecommand{\Oij}{\(\mathrm{O};\vec{\imath},\vec{\jmath}\)}
\providecommand{\Oijk}{\(\mathrm{O};\vec{\imath},\vec{\jmath},\vec{k}\)}
\providecommand{\Ouv}{\(\mathrm{O};\vec{u},\vec{v}\)}
\providecommand{\intervFF}[2]{\ensuremath{[#1;#2]}}
\renewcommand{\euro}{\ifmmode\text{\texteuro}\else\texteuro\fi}
\DeclareUnicodeCharacter{00B0}{\ifmmode^{\circ}\else\textdegree\fi}
\DeclareUnicodeCharacter{20AC}{\ifmmode\text{\texteuro}\else\texteuro\fi}
% Énoncés : Éducation nationale. Transcription et corrigés : APMEP, Denis Vergès et contributeurs. Découpage et classement : CoopMaths, Sébastien Lozano, Éric Elter et contributeurs. Adaptation de la présentation : Maths Explorer.
% CC BY-SA : https://coopmaths.fr/www/about/
% https://coopmaths.fr/alea/static/bac/2022/tex/bac_2022_10_sujet2_caledonie_2.tex
\begin{document}
\noindent\textit{APMEP / CoopMaths. CC BY-SA. Présentation adaptée par Maths Explorer.}\par\medskip
\textbf{Principaux domaines abordés :} fonctions, fonction logarithme.

\medskip

Soit $f$ la fonction définie sur l'intervalle $]0~;~+\infty[$ par :

\[f(x) = x \ln (x) - x - 2.\]

On admet que la fonction $f$ est deux fois dérivable sur $]0~;~+\infty[$.

On note $f'$ sa dérivée, $f''$ sa dérivée seconde et $\mathcal{C}_f$ sa courbe représentative dans un repère.

\medskip

\begin{enumerate}
\item 
	\begin{enumerate}
		\item Démontrer que, pour tout $x$ appartenant à $]0~;~+\infty[$, on a $f'(x) = \ln (x)$.
		\item Déterminer une équation de la tangente $T$ à la courbe $\mathcal{C}_f$ au point
d'abscisse $x = $e.
		\item Justifier que la fonction $f$ est convexe sur l'intervalle $]0~;~+\infty[$.
		\item En déduire la position relative de la courbe $\mathcal{C}_f$ et de la tangente $T$.
	\end{enumerate}	
\item 
	\begin{enumerate}
		\item Calculer la limite de la fonction $f$ en $0$.
		\item Démontrer que la limite de la fonction $f$ en $+\infty$ est égale à $+\infty$.
	\end{enumerate}
\item Dresser le tableau de variations de la fonction $f$ sur l'intervalle $]0~;~+\infty[$.
\item 
	\begin{enumerate}
		\item Démontrer que l'équation $f(x) = 0$ admet une unique solution dans
l'intervalle $]0~;~+\infty[$. On note $\alpha$ cette solution.
		\item Justifier que le réel $\alpha$ appartient à l'intervalle ]4,3~;~4,4[.
		\item En déduire le signe de la fonction $f$ sur l'intervalle $]0~;~+\infty[$.
	\end{enumerate}	
\item On considère la fonction \texttt{seuil} suivante écrite dans le langage Python :

On rappelle que la fonction \texttt{log} du module \texttt{math} (que l'on suppose importé)
désigne la fonction logarithme népérien ln.

\begin{center}
\begin{tabular}{|l l|} \hline
\texttt{def}& \texttt{seuil(pas) :}\\
&\texttt{x=4.3}\\
&\texttt{while x*log (x) - x - 2 < 0:}\\
&\quad \texttt{x=x+pas}\\
&\texttt{return x}\\ \hline
\end{tabular}
\end{center}

Quelle est la valeur renvoyée à l'appel de la fonction \texttt{seuil(0.01)} ?

Interpréter ce résultat dans le contexte de l'exercice.
\end{enumerate}

\bigskip

\end{document}
